\section{Introduction}\label{sec:introduction}

Chromatic height organizes both the deformation theory of formal groups
and the tensor structure of stable homotopy theory.  Two classification
problems ask whether height accounts for all the relevant closed
subobjects: the invariant ideals in a Lubin--Tate deformation ring, and
the thick tensor ideals of dualizable Morava-local spectra.  We prove
both classifications, using the arithmetic problem as the main step.

Fix a prime \(p\) and an integer \(n\geq1\).  Let \(\Gamma_n\) be the
height-\(n\) Honda formal group, let \(A_n\) be its endomorphism ring over
\(\overline{\mathbb F}_p\), and put
\[
 E_0=W(\mathbb F_{p^n})[[u_1,\ldots,u_{n-1}]],
 \qquad I_0=0,\qquad
 I_h=(p,u_1,\ldots,u_{h-1})\quad(1\leq h\leq n).
\]
The ring \(E_0\) represents its marked deformations.  The ideals \(I_h\)
are the height ideals: in characteristic \(p\), their vanishing specifies
where the formal group has height at least \(h\).
The nonextended Morava stabilizer \(A_n^\times\) acts by changing the
marking.  The canonical extended stabilizer \(G_n\) also includes the
residue-field Galois action.  The unit ideal is listed separately from
the proper ideals.

\begin{theorem}[Stabilizer-invariant ideals]\label{thm:invariant-radicals}
Let \(U\subseteq A_n^\times\) be any open subgroup.
The proper prime ideals of \(E_0\) stabilized by \(U\) are precisely
\[
 0,\ I_1,\ldots,I_n.
\]
The \(U\)-invariant radical ideals are precisely
\[
 0,\ I_1,\ldots,I_n,\ E_0.
\]
The latter list also gives the radical ideals invariant under the full
canonical \(G_n\)-action.
\end{theorem}

This proves the invariant-ideal assertion called Chai's Hope in the
finite-residue-field formulation of
\cite[Hope~4.1 and Remark~4.2]{BHN22}, and extends that formulation
from the full stabilizer to every open subgroup.  The open-subgroup
scope also occurs in Chai's original closed-fiber question, as explained
in \cite[Remark~4.2]{BHN22}.  The second classification concerns Honda
Morava \(E\)-theory \(E\) with coefficient ring \(E_0\), and the
homotopy category \(\mathcal D_{(n,p)}\) of dualizable \(K(n)\)-local
spectra.  Its tensor product and unit are
\[
 X\otimes Y=L_{K(n)}(X\wedge Y),
 \qquad \mathbf 1=L_{K(n)}S.
\]
A thick tensor ideal is a full subcategory closed under equivalences,
finite sums, suspensions and desuspensions, cofiber sequences, retracts,
and tensoring with every object of \(\mathcal D_{(n,p)}\).
Dualizability is essential here: it is more general than compactness in
the \(K(n)\)-local category \cite[Theorems~8.5--8.6]{HS99}.

A finite \(p\)-local spectrum \(F\) has \emph{type \(k\)} if
\(K(i)_*F=0\) for \(i<k\) and \(K(k)_*F\neq0\), with
\(K(0)=H\mathbb Q\).  Such spectra exist in every type: start with the
sphere and the degree-\(p\) cofiber, and then take cofibers of the
self-maps supplied by the Hopkins--Smith periodicity theorem
\cite[Definition~8 and Theorem~9]{HopkinsSmith98}.
For \(0\leq k\leq n\), choose a finite type-\(k\) spectrum \(F(k)\)
and put
\[
 \mathcal D_k=\big\langle L_{K(n)}F(k)\big\rangle_\otimes,
 \qquad \mathcal D_{n+1}=0,
\]
where the brackets mean the generated thick tensor ideal in
\(\mathcal D_{(n,p)}\).

\begin{theorem}[Hovey--Strickland conjecture]\label{thm:main}
For every prime \(p\) and every \(n\geq 1\), the thick tensor ideals of
\(\mathcal D_{(n,p)}\) are exactly
\[
 \mathcal D_{(n,p)}=\mathcal D_0
 \supsetneq\mathcal D_1\supsetneq\cdots
 \supsetneq\mathcal D_n\supsetneq\mathcal D_{n+1}=0.
\]
In particular, there are \(n+2\) such ideals, and each \(\mathcal D_k\)
is independent of the chosen finite type-\(k\) spectrum.
\end{theorem}

The Balmer spectrum of a tensor category is the space of its prime thick
tensor ideals.  Corollary~\ref{fin:balmer} identifies this spectrum here
with the \(n+1\) proper members of the displayed chain, including the
zero ideal.  It also specifies their specialization order and the
support of each finite type generator.

\subsection{History and the two classification problems}

The thick subcategory theorem of Hopkins and Smith classifies thick
subcategories of finite spectra by chromatic type, while their
periodicity theorem supplies the corresponding self-maps
\cite[Theorems~7 and~9]{HopkinsSmith98}.  Hovey and Strickland proposed
the dualizable \(K(n)\)-local counterpart in their study of Morava
\(K\)-theory and localization \cite[Section~12, p.~61]{HS99}.
Their Problem~16.8 also asks directly about stabilizer-invariant primes
and radical ideals of the deformation ring; it explains why the
invariant-prime theorem for ordinary complex cobordism does not by
itself answer this question.

Chai studied the action on the closed fiber of Lubin--Tate space and
proposed that invariant irreducible formal loci should be height loci
\cite{Chai96}.  We use the precise formulation and comparison of
coefficient fields and group actions in
\cite[Section~4, Hope~4.1 and Remark~4.2]{BHN22}.
Formal rigidity also has an established role in Chai's work on
\(p\)-divisible formal groups \cite{Chai08}.  The formal-closure
argument below belongs to this general circle of ideas, but starts
from a subgroup of Honda-group points and proves its required
annihilator statement directly.

Barthel, Heard, and Naumann proved the Hovey--Strickland conjecture at
height two for every prime and established the general implication from
Chai's invariant-ideal assertion to the spectral classification
\cite[Theorems~4.9 and~4.15]{BHN22}.  The descent and support methods in
that implication build on Mathew's work on Morava \(E\)-theory descent
\cite{Mathew16} and Balmer's tensor-nilpotence criterion for surjectivity
on spectra \cite{Balmer18}.  We recall the exact support comparison in
Proposition~\ref{fin:support}.  Its direction is unrestricted in \(p\).
The inequality \(2p-2>n^2+n\) in \cite[Theorem~4.13]{BHN22}
concerns the converse, which uses realization of Morava modules.

The generic-fiber part has earlier Lie-theoretic precedents.
Barthel--Heard--Naumann credit Chai with an unpublished
characteristic-zero invariant-ideal result and use the Gross--Hopkins period map to study
projective stabilizer orbits
\cite[Remark~4.2 and Proposition~4.18]{BHN22}.
We retain a local analytic-germ proof of the needed conclusion.
The main issue addressed here is the characteristic-\(p\) height
strata, where infinitesimal stabilizer directions alone do not give
the required control of invariant equations.

Recent generation and comparison results concern related categories.
Lee and Pstr{\k a}gowski identify the category of locally fp spectra,
which contains the dualizable \(K(n)\)-local spectra as well as
nondualizable objects
\cite[Example~1.6 and Theorems~1.7--1.9]{LP26}.
Surjectivity of the map from the \(K(n)\)-local to the \(E(n)\)-local
Balmer spectrum was established by Barthel--Heard--Naumann
\cite[Remark~3.6]{BHN22}; Balmer--Sanders give a later general treatment
\cite[Proposition~8.11 and Remark~8.13]{BalmerSanders}.
These results provide generation and surjectivity statements.
The arithmetic theorem above classifies the invariant supports needed
to determine the entire ideal lattice.

\begin{figure}[htbp]
\centering
\begin{tikzpicture}[
  box/.style={draw,rounded corners=2pt,align=center,font=\small,
              inner sep=5pt,text width=5.6cm},
  merge/.style={box,text width=9.6cm},
  flow/.style={-{Latex[length=2mm]},thick}
]
\node[box,minimum height=1.8cm] (charp) at (-3.35,0)
  {Characteristic $p$\\Formal closure\\Bundle nondegeneracy\\
   Tate coordinates (Sections~\ref{sec:formal}--\ref{sec:deformations})};
\node[box,minimum height=1.8cm] (charzero) at (3.35,0)
  {Characteristic zero\\Gross--Hopkins period map\\
   Projective directions (Section~\ref{sec:classification})};
\node[box,below=5mm of charp] (density)
  {Local orbit density (Section~\ref{sec:orbits})};
\node[merge,below=8mm of density.south,xshift=3.35cm] (primes)
  {Open-subgroup invariant primes and radical ideals};
\node[merge,below=5mm of primes] (support)
  {All-prime support comparison of Barthel--Heard--Naumann};
\node[merge,below=5mm of support] (spectra)
  {$n+2$ thick tensor ideals and $n+1$ Balmer points};
\draw[flow] (charp) -- (density);
\draw[flow] (density.south) -- (density.south |- primes.north);
\draw[flow] (charzero.south) -- (charzero.south |- primes.north);
\draw[flow] (primes) -- (support);
\draw[flow] (support) -- (spectra);
\end{tikzpicture}

\caption{The two orbit arguments classify invariant primes.  Only the
characteristic-\(p\) branch uses the formal, bundle, and deformation
constructions of Sections~\ref{sec:formal}--\ref{sec:deformations}.
The arithmetic branches join before the all-prime support comparison.}
\label{fig:proof-flow}
\end{figure}

\subsection{The argument}

Set \(C=\Cp^\flat\) and \(k=\overline{\mathbb F}_p\), viewed as the
coefficient field in \(C\), and write
\(\mm_C=\{c\in C:|c|<1\}\).  Fix a characteristic-\(p\) point \(x\)
of formal height \(1\leq h<n\).  The main task is to show that a
nonzero analytic germ at \(x\) on the height-at-least-\(h\) parameter locus
\(I_h=0\) cannot vanish on the stabilizer orbit of \(x\).
The finite kernels of the universal deformation give a height-\(n\)
algebraic \(p\)-divisible group.  At \(x\), its connected part has
height \(h\) and its \'{e}tale quotient has height \(r=n-h\).
The universal-cover comparison lets us regard a basis
\(z_1,\ldots,z_r\) of its \'{e}tale Tate module as points of the
perfected height-\(n\) Honda group.  Specializing these fixed vectors
at nearby parameters in this locus gives analytic functions \(\theta_i\) vanishing
at \(x\).  Their vanishing tests whether the vectors remain in the
nearby Tate lattice.  Section~\ref{sec:deformations} proves that their
parameter differential is invertible: a vector in its kernel splits
the \'{e}tale quotient over the dual numbers, fixed-height rigidity
makes the connected factor constant, and the Kodaira--Spencer map
then kills the parameter vector
(Proposition~\ref{def:coordinates}).

Put \(q=p^h\).  For a fixed \(d\in A_n\), let
\(y_N=(1-p^Nd)\cdot x\); these points tend to \(x\).
Stabilizer transport relates \(\theta_i(y_N)\) to multiplication by
\(p^N\) in the formal group at \(y_N\).  Integral comparison series
intertwining multiplication by \(p\) with that of the height-\(h\)
Honda group turn this relation into coordinatewise Frobenius powers.
Applying these series to the specializations of \(z_i\) and
\(dz_i\) at \(y_N\) gives tuples \(U_N\) and \(s_N\), respectively,
in \(\mm_C^r\).  Theorem~\ref{orb:density} proves the exact identity
and limit
\[
 U_N=s_N^{q^N},\qquad
 s_N\longrightarrow s(d)=\bigl(\nu_x(dz_1),\ldots,\nu_x(dz_r)\bigr),
\]
where \(\nu_x\) is specialization at \(x\) followed by an integral
formal identification with the height-\(h\) Honda group.
Proposition~\ref{def:natural-comparison} makes \(\nu_x\) a nonzero
natural map on perfectoid test spaces.  The corresponding change of
parameter coordinates is formal.  To compare it with the numerical
tuples \(U_N\), we match any fixed finite number of analytic
coefficient germs and bound the remaining integral tails
(Lemma~\ref{orb:finite-jets}).  The coefficient germs may have
shrinking domains; the estimates require only one finite jet at a time.

Suppose a nonzero analytic germ \(P\) vanishes on the orbit.  Its expansion
in the formal coordinates has fixed coefficients \(B_\ell\), chosen
before \(d\).  Finite truncation followed by \(q^N\)-th roots gives
bounds on the finite sums
\(\sum_{|\ell|\leq m} B_\ell^{1/q^N}s(d)^\ell\).
Frobenius separation (Lemma~\ref{formal:separation}) turns these
bounds into a single nonzero formal equation over \(k\) on the subgroup
\[
 S_x=\{\bigl(\nu_x(dz_1),\ldots,\nu_x(dz_r)\bigr):d\in A_n\}.
\]
The formal-closure argument in Section~\ref{sec:formal} then produces
a nonzero row of Honda endomorphisms annihilating \(S_x\)
(Proposition~\ref{formal:closure}).

The contradiction is geometric.  The Fargues--Fontaine curve
\cite{FF18} and Le Bras's theorem on Banach--Colmez sheaves
\cite{LeBras18} turn this natural endomorphism relation into a
vector-bundle relation.  A degree-one determinant argument promotes
the \(\Qp\)-independence of \(z_1,\ldots,z_r\) to independence at the
generic point of the curve (Lemma~\ref{bundles:independence}).
There the division algebra action spans the full matrix algebra.
A matrix selecting one of these independent vectors contradicts the
nonzero endomorphism row
(Proposition~\ref{bundles:nondegeneracy}).  Thus \(P\) must be zero.
This vanishing property is the \emph{local orbit density} proved in
Theorem~\ref{orb:density}.

Section~\ref{sec:classification} uses this result to classify invariant
primes containing \(p\).  Primes in characteristic zero are handled
separately by the Gross--Hopkins period map \cite{GH94,HG94Bundles}.
A finite minimal-prime argument classifies invariant radical ideals,
and the established support comparison transfers the height chain to
thick tensor ideals and their Balmer spectrum.  The final description
also identifies the support map from \(\Spec(E_0)\): each height stratum
maps to one Balmer point.  Figure~\ref{fig:proof-flow} records the two
arithmetic branches and where they enter the spectral argument.
